\chapter{Admissibility Preservation}
\label{ch:preservation}

\section{Statement}

The question addressed in this chapter is what it means for a
trajectory through admissibility structure to remain coherent. The
answer given here is that coherence is not an additional condition
imposed on top of admissibility; it is what nontriviality of the
admissibility relation amounts to when examined along a trajectory
rather than at a single state.

\begin{definition}[Trajectory]
A trajectory is a sequence of states $S_0, S_1, \dots, S_t$ such that
$S_i \adm{R} S_{i+1}$ for each $i < t$, where $R$ may itself vary with
$i$ if the underlying admissibility structure is time-indexed as $R_i$.
\end{definition}

\begin{definition}[Degeneracy]
A trajectory degenerates at step $i$ if $\reach{S_i}$ is empty (no
continuation is admissible) or a singleton entirely determined by
$S_0, \dots, S_{i-1}$ (exactly one continuation is admissible, and it
was already fixed before $S_i$ was reached).
\end{definition}

\begin{theorem}[Admissibility Preservation]
\label{thm:preservation}
A trajectory $S_0, \dots, S_t$ is coherent if and only if it does not
degenerate at any step $i \le t$; equivalently, if and only if
$|\reach{S_i}|$ exceeds a nontrivial floor relative to what has already
been committed by $S_0, \dots, S_{i-1}$, at every step.
\end{theorem}

\begin{proof}
Coherence, as the term is used throughout this program, has no content
apart from the admissibility relation itself: to say a trajectory is
coherent is to say that at each step, the relation $x \adm{R} y$
continues to discriminate admissible from inadmissible continuations
non-vacuously. Suppose the trajectory degenerates at step $i$ because
$\reach{S_i} = \varnothing$. Then no $y$ satisfies $S_i \adm{R} y$, and
the relation, restricted to continuations from $S_i$, is uniformly
false; there is no sense in which the trajectory can be said to continue
admissibly past $S_i$, coherently or otherwise, since no continuation is
admissible at all. Suppose instead the trajectory degenerates because
$\reach{S_i}$ is a singleton already fixed by $S_0, \dots, S_{i-1}$.
Then the relation restricted to continuations from $S_i$ is uniformly
true of exactly one candidate and uniformly false of every other; the
admissibility relation is doing no discriminating work at step $i$, since
its outcome was determined before $S_i$ was reached, and the statement
"$S_i \adm{R} S_{i+1}$ is coherent" becomes true independently of
anything about $S_i$ itself. In either degenerate case, the relation
$\adm{R}$ fails to be a nontrivial constraint at step $i$: it is either
vacuously unsatisfiable or vacuously predetermined. Conversely, if
$\reach{S_i}$ exceeds a nontrivial floor at every step --- if there is
always more than one admissible continuation not already fixed by the
trajectory's own history --- then the admissibility relation is doing
genuine discriminating work at every step, which is what coherence
consists in. This establishes the biconditional.
\end{proof}

\section{Discussion}

The proof above is deliberately minimal, and its minimality is the
point. Coherence has sometimes been treated, informally, as though it
were an additional property a trajectory could have or lack independent
of the admissibility relation that generates it --- as though a
trajectory could be admissible at every step and yet somehow
incoherent, or coherent despite violating admissibility. Theorem
\ref{thm:preservation} rules this out: coherence and non-degeneracy of
the admissibility relation are the same property under two names. This
matters directly for the definition of monotonic continuation developed
elsewhere in this program, where the ordering $S_t \le S_{t+1}$
presupposes, without stating, that $\reach{S_t}$ remains nontrivial;
Theorem \ref{thm:preservation} supplies the missing lemma that makes
this presupposition explicit and shows exactly what fails when it does
not hold.

\begin{corollary}
\label{cor:preservation-minimum-support}
Any operational test for coherence --- such as a minimum support ratio
required of a discourse trajectory before it is judged coherent --- is
correctly understood as an operational proxy for $|\reach{S_i}|$
remaining above a nontrivial floor, not as an independent criterion.
Where such a proxy diverges from the underlying reachable-set condition,
the proxy, not the theorem, should be revised.
\end{corollary}

Corollary \ref{cor:preservation-minimum-support} is included because
domain-specific coherence tests --- support ratios, consistency scores,
and similar quantities --- are frequently introduced as though they were
primitive definitions of coherence in their own right. Theorem
\ref{thm:preservation} reduces the space of legitimate such tests to
those that track, however indirectly, the nondegeneracy of the
admissibility relation, and gives a principled way to evaluate a
proposed test: ask whether it can be false while $\reach{S_i}$ remains
nontrivial, or true while $\reach{S_i}$ has degenerated. Either case
indicates the proxy has diverged from what coherence actually
consists in.
